\section{Construções em Fin}

A categoria \textbf{Fin} é completa e cocompleta, o que permite a existência de todos os limites e colimites finitos. A implementação em MATLAB utiliza predominantemente operações sobre vetores e matrizes de índices.

\subsection{Fundamentação Teórica}
\begin{itemize}
    \item \textbf{Produto ($A \times B$):} O objeto resultante tem cardinalidade $n \times m$. As projeções são geradas para cobrir todas as combinações do produto cartesiano.
    \item \textbf{Coproduto ($A \sqcup B$):} União disjunta com objeto $n + m$. As injeções mapeiam o primeiro conjunto para o início do novo intervalo e o segundo para o restante.
    \item \textbf{Equalizador:} O subconjunto do domínio onde $f(i) = g(i)$.
    \item \textbf{Coequalizador:} Espaço quociente da relação $f(i) \sim g(i)$. É modelado computacionalmente através de componentes conexas de um grafo onde as arestas ligam as imagens dos morfismos.
    \item \textbf{Pullback:} O conjunto de pares $(x, y)$ que satisfazem a equação $f(x) = g(y)$.
\end{itemize}

\subsection{Implementação Computacional}
\begin{lstlisting}[caption={Construções Universais em Fin}]
% Produto e Projeções
[P1, P2] = ndgrid(1:n, 1:m);
p_A = P1(:)'; p_B = P2(:)'; 

% Equalizador
eq_map = find(f == g);
eq_obj = length(eq_map);

% Pullback (f: A -> C, g: B -> C)
[X, Y] = ndgrid(1:length(dom_f), 1:length(dom_g));
pb_indices = find(f(X) == g(Y));

% Coequalizador (Esquema de componentes conexas)
% Requer a criacao de um grafo com arestas (f(i), g(i))
% e calculo de 'conncomp' para identificar as classes.
\end{lstlisting}